% ============================================================
% PREAMBLE -- Rapport de stage PFA1 (format UM6P)
% Police Carlito (clone metrique libre de Calibri), taille 12.
% Compiler avec XeLaTeX / tectonic.
% ============================================================

\usepackage[french]{babel}
\usepackage[margin=2.4cm, top=2.6cm, bottom=2.4cm]{geometry}

% ------------------------------------------------------------
% POLICE : Carlito = clone metrique libre de Calibri
% ------------------------------------------------------------
\usepackage{fontspec}
\setmainfont{Carlito}[
  Extension      = .ttf,
  Path           = /usr/share/fonts/carlito/,
  UprightFont    = Carlito-Regular,
  BoldFont       = Carlito-Bold,
  ItalicFont     = Carlito-Italic,
  BoldItalicFont = Carlito-BoldItalic
]
\setmonofont{DejaVu Sans Mono}[Scale=0.85]

\usepackage{amsmath}
\usepackage{amssymb}

\usepackage{graphicx}
\graphicspath{{../plots/}{./figures/}}

\usepackage{booktabs}
\usepackage{array}
\usepackage{tabularx}
\usepackage{adjustbox}
\usepackage{enumitem}
\usepackage[table]{xcolor}
\usepackage{caption}
\usepackage{float}

\usepackage{tikz}
\usetikzlibrary{arrows.meta, positioning}

\usepackage{fancyhdr}
\usepackage{titlesec}
\usepackage{titletoc}

\usepackage[hidelinks]{hyperref}

% Intitules conformes aux consignes de l'ecole
\addto\captionsfrench{%
  \renewcommand{\listfigurename}{Liste des figures}%
  \renewcommand{\listtablename}{Liste des tableaux}%
  \renewcommand{\bibname}{Références bibliographiques}%
}

% ------------------------------------------------------------
% PALETTE
% ------------------------------------------------------------
\definecolor{accent}{RGB}{17,112,120}      % teal -- titres
\definecolor{caplabel}{RGB}{26,105,150}    % bleu -- etiquettes figures/tables
\definecolor{umorange}{RGB}{215,74,42}     % orange UM6P -- filets
\definecolor{rulegray}{RGB}{120,120,120}

% ------------------------------------------------------------
% LEGENDES : etiquette coloree en gras, texte justifie
% ------------------------------------------------------------
\captionsetup{
  font=small,
  labelfont={bf,color=caplabel},
  labelsep=period,
  justification=justified,
  singlelinecheck=false,
  skip=6pt
}
\captionsetup[figure]{justification=justified}
\captionsetup[table]{justification=justified}

\renewcommand{\arraystretch}{1.15}
\newcolumntype{Y}{>{\centering\arraybackslash}X}
\newcommand{\rowrule}{\midrule[0.2pt]}

% ------------------------------------------------------------
% EN-TETES / PIEDS
% ------------------------------------------------------------
% Chaque style est defini par \fancypagestyle et re-execute son propre corps
% a l'activation : sans cela, le \fancyhf{} du style "plain" efface
% definitivement les en-tetes du style principal (registres globaux fancyhdr).
\setlength{\headheight}{15pt}

\fancypagestyle{main}{%
  \fancyhf{}%
  \renewcommand{\headrulewidth}{0.4pt}%
  \renewcommand{\footrulewidth}{0pt}%
  \fancyhead[L]{\small\nouppercase{\leftmark}}%
  \fancyhead[R]{\small Rapport de stage PFA1}%
  \fancyfoot[C]{\thepage}%
}

\fancypagestyle{plain}{%
  \fancyhf{}%
  \renewcommand{\headrulewidth}{0pt}%
  \renewcommand{\footrulewidth}{0pt}%
  \fancyfoot[C]{\thepage}%
}

\pagestyle{main}

% ------------------------------------------------------------
% TITRES
% ------------------------------------------------------------
\titleformat{\chapter}[display]
  {\normalfont\huge\bfseries\color{accent}}
  {\large\chaptertitlename\ \thechapter}
  {8pt}
  {\Huge}
  [{\color{umorange}\titlerule[1.2pt]}]
\titlespacing*{\chapter}{0pt}{-28pt}{16pt}

\titleformat{\section}
  {\large\bfseries\color{accent}}
  {\thesection}{0.7em}{}
\titlespacing*{\section}{0pt}{2.8ex plus 1ex minus .2ex}{1.2ex plus .2ex}

\titleformat{\subsection}
  {\normalsize\bfseries\color{accent}}
  {\thesubsection}{0.7em}{}
\titlespacing*{\subsection}{0pt}{2.2ex plus 1ex minus .2ex}{0.9ex plus .2ex}

\titleformat{\subsubsection}
  {\normalsize\bfseries}
  {\thesubsubsection}{0.7em}{}
\titlespacing*{\subsubsection}{0pt}{1.8ex plus 1ex minus .2ex}{0.7ex plus .2ex}

% ------------------------------------------------------------
% STYLE TIKZ partage par les schemas
% ------------------------------------------------------------
\tikzset{
  rbox/.style={draw, rounded corners, fill=accent!10, align=center,
               minimum height=1.0cm, text width=#1, font=\small},
  rbox/.default=2.6cm,
  flow/.style={-{Latex[length=2.2mm]}, thick}
}

% ------------------------------------------------------------
% DIVERS
% ------------------------------------------------------------
\setlength{\parskip}{0.42em}
\setlength{\parindent}{0pt}
\setlist{itemsep=2pt, topsep=3pt, parsep=0pt}
